\section{Problem description}

The goal of the project was thus to derive a map of possible optimal designs of the influent diet mixture, of which one is to be selected by plant operators depending on the "risk level" they are willing to take in face of uncertainties.
To do so, the methodology applied considered the robustification of a linear programming (LP) problem.

The model considered a biomethane plant in operation as an "intermediate conversion" layer between an \emph{upstream} market of waste/byproducts (multiple inputs) and the \emph{downstream} market of energy retailing (single output).
It focuses on the optimization of the \emph{upstream} side of the plant reactor i.e. finding the influent diet that results in the highest biomethane production with the lowest operational expenditure, given the time-varying market, the geographical and the legislative boundary conditions. 

The goal was thus to have a tool that decision-makers can use to modify/update the influent diet for the future operational time window, according to the present market conditions.
Decisions have thus to be taken in terms of 'how much' of 'which substrate' to be purchased from 'which seller' available in the nearby territory.
The problem is solved based on data that can be made available to plant managers. The spòutions suggest on which are the optimal contracts to be negotiated based on the dataset of recorded past realizations of substrates price and availability.

Given the time dependency of market conditions (ie. availability and purchase cost of substrates), a 'puchase horizon'/'operating window' can be arbitrarily set by the decision maker (usually between 3-12 months to guarantee a safe and "stable" operation).
Substrates are usually stored at the plant for long 'operating window' (es. to buffer seasonal availability of different substrates such as crops and by-products). Nevertheless, the delivery of the purchased quantities can also be scheduled in time in accordance with the seller. Also, long-time stocks are not suggested as they may introduce a noteworthy degradation of the substrate performances in the reactor.

The plant is a priori authorized to transform only a subset of item types available on the market. From the legislative point of view, the case-study that was taken into account focused on an 'agricultural' plant, in which biomethane has to be directly injected to the grid for third use (i.e. no wastes in the influent diet are allowed to have a guaranteed access to the incentives).
The case-study refers to a plant located in the Cremona province, Northern Italy.

%%%%%%%%%%%%%%%%%%%%%%%%%%%%%%%%%%%%%%%%%%%%%%%%%%%%%%%%%%%%%%%%%%%%%%%%%%%%%%%%
\subsection{Objective function}
The objective function was the maximisation of the profit from the selling of the output that can be produced with the purchased co-substrates and to be fed to the reactor until the next 'purchase horizon'. The profit was computed as the different between the revenues from biomethane selling and the costs related to the purchase and transportation of substrates. 
Substrate prices are usually referred to the overall weight unit of such substrate (i.e. ton of fresh matter (FM)). Nevertheless, each substrate is composed by a fraction of water and "total solid" (TS) content. The latter is also usually partitioned in the "ash" and "volatile solid" (VS) contents. The VS of each substrate is characterized by its content of organic acids, biodegradable carbohydrates, proteins and lipids, from which anaerobic degradation results in the substrate-specific biochemical methane potential (BMP (NL$_{ch4}$kg$_{VS}^{-1}$)) i.e. the substrate-specific conversion yield.

In a \underline{first simplified approach}, the BMP of each substrate was considered equal to the maximum conversion yield i.e. the conversion achieved at the end of lab-scale batch-tests on that specific substrate (BMP$_\infty$), that usually last between 30-50 days to guarantee complete degradation of the biodegradable fraction. The BMP$_{\infty}$ is thus the asymptotic value of the cumulative amount of methane produced out of a certain amount of $VS$ and can be usually found in literature.

Nevertheless, in continuous or semi-continuous reactor feeding (typical of full-scale configurations), the reactor is characterized by a certain expected hydraulic residence time (HRT) of the influent organic matter, which is usually lower with respect to the duration of lab-scale batch-tests, in order to guarantee a higher biomethane flowrate. Thus, the actual operation is less efficient with respect to the maximum extent of conversion, for a higher productivity. Indeed, each substrate is not only characterized by its BMP$_\infty$ value, but also by the time-dependent BMP curve, which maps the conversion yield as a function of time from the batch-test starting point (i.e. similar to the residence time in a full-scale reactor). HRT is one of the main process parameter and its optimization was carried out along with diet optimization, in a \underline{more realistic approach}, slightly modifying the problem formulation: since the obtained objective function was non-linear in the decision variables, one LP problem was solved for each HRT.
The underlying assumption of 'additivity' between the BMP of the different co-substrates at a given HRT is still a simplification, as inhibition, competitive or cooperative effects may arise; from this consideration, the obtained plant productivity is indeed an upper bound to real in-plant operation and optimal diet must to further validated with the ADM1 model.
Eventually, it is noteworthy to mention that the value of the objective function that is obtained is not representative for the true net profit, since it is neglecting other operational expenditures such as energy self-consumption, etc. were neglected for simplicity.
%%%%%%%%%%%%%%%%%%%%%%%%%%%%%%%%%%%%%%%%%%%%%%%%%%%%%%%%%%%%%%%%%%%%%%%%%%%%%%%%
\subsection{Constraints}
\begin{itemize}
\item \emph{Seller substrates availability}. The amount of each substrate that can be purchased from each seller must be upper bounded by its availability. Sellers are either 'agribusinesses' or 'industries'. 
\item \emph{Regulating authority constraints}. To warrant access to incentives, the regulating autorithy enforced some bound constraints on the diet composition. In particular, the definitions of 'product', 'by-product', 'animal slurry' and 'second harvest crop' are given to cluster substrates and enforcing percentage bounds on the clustered diet composition.
\item \emph{Hydraulic retention time}. The "processable" total flowrate of the influent co-substrate mixture shall be limited with bound constraints set by expert knowledge on process stability, and limits imposed by the available stock capacity of digestate and/or by digestate land spreading legislative limits (i.e. if digestate can't be valorized, it has to spread over a certain amount of available hectares of land, high enough to guarantee a maximum allowed nitrogen load per hectare).
\item \emph{Technical constraint}. A set of K technical limitations defined by heuristic knowledge shall be applied to find a co-digestion diet that shall, at least preliminarily, be able to lead to stable operations. They are based on typically available substrate characterization and are applied in the form of bound constraints over the weighted average of the \emph{k-th} characteristic of the co-substrates mixture.
\item \emph{Residual digestate BMP}. Imposing a limit to the residual BMP ($BMP_{res,max}$) in the digestate effluent is intended to have the effect to mold the feasibility region toward higher HRTs, with a consequent lower risk for reactor instability (i.e. acidification) that the LP model is not able to grasp at low HRTs (in contrary to the ADM1 model). Upcoming limitation by the regulating authority is also foreseen in this sense, as it has the effect to limit the residual load of undegradable organic matter in the effluent (potential pollution of the environment by digestate land spreading) too. This constraint has meaning only for the "realistic approach" mentioned above.
\item \emph{Non-negativity constraint}. The plant is not authorized to sell substrates to any third party.
\end{itemize}
%%%%%%%%%%%%%%%%%%%%%%%%%%%%%%%%%%%%%%%%%%%%%%%%%%%%%%%%%%%%%%%%%%%%%%%%%%%%%%%%
\subsection{Uncertainty sources}
When dealing with organic substrates and unsaturated markets, a lot of uncertainty comes into play.
For semplicity and demostrative purposes, only the following sources of uncertainty were considered:
\begin{itemize}
\item \emph{Substrate purchase price}. The purchase price is the main operational cost source for plant managers (along with maintenance, that was not taken into consideration in the model). Before a contract is negotiated, each substrate price is influenced by its primary market equilibrium conditions, which volatility depends on uncountable factors and on the substrate availability itself. In the case of industrial by-products, the market is not saturated, so that the conditions are even more uncertain. For animal slurries, the price is in general dictated by the nutrient content, as for their potential valorization in the fertilizers market; in this sense, the uncertainty in the purchase price of slurries has not to be attributed to a variable nutrient content (which was hypothesized as fixed), but to the volatility of the market price of nitrogen (N), phosphorous (P) and potassium (K). It entails a time-dependent (primarily seasonal) dynamic.

\item \emph{Substrate availability}. The uncertainty in substrate availability is mainly driven by seasonal variations and weather events. In the case of industrial by-products, the availability depends on the production volume of the sellers for their own primary market. It entails a time-dependent dynamic (es. some crops are cultivated and thus more available in summer, while others in winter).

\item \emph{BMP$_{\infty}$}. The values that can be found in literature are characterised by a certain variance because of the degree of variability in the testing conditions and the different "specific substrates" (with different technical characteristics) that are lumped under the broader 'substrate' umbrella (es. different cultivar of maize silage or maize silage growth in different times of the year are lumped within the 'maize silage' indexing). In case a more detailed database would be available to plant managers, substrate indexing could be de-lumped and min-max variability reduced.
\end{itemize}

On the other hand, the uncertainty related to substrate technical characterization (that in general may also present some non-linear correlation to the substrate $BMP_{\infty}$) was neglected for simplicity; this is also consistent with the bounds of the technical constraints being uncertain/dictated by expert knowledge too. Indeed, technical constraints are meant to give a first preliminary 'feasibility' check of the purchased diet, but further validation with the ADM1 model is for sure required for prediction of actual in-plant operation. 
